% Generated draft: requires mathematical and source review.
\documentclass[11pt]{article}
\usepackage[margin=25mm]{geometry}
\usepackage{fontspec}
\setmainfont{Latin Modern Roman}
\usepackage{amsmath,amssymb,mathtools,unicode-math}
\setmathfont{Latin Modern Math}
\usepackage{enumitem,needspace,xurl,hyperref}
\setlength{\parindent}{0pt}
\setlength{\parskip}{5pt}
\setlength{\emergencystretch}{3em}
\newcommand{\N}{\mathbb{N}}
\newcommand{\Z}{\mathbb{Z}}
\newcommand{\Q}{\mathbb{Q}}
\newcommand{\R}{\mathbb{R}}
\newcommand{\C}{\mathbb{C}}
\newcommand{\F}{\mathbb{F}}
\newcommand{\A}{\mathbb{A}}
\newcommand{\PP}{\mathbb{P}}
\newcommand{\E}{\mathbb{E}}
\newcommand{\eps}{\varepsilon}
\newcommand{\dd}{\,\mathrm{d}}
\begin{document}
\section*{UnsolvedMath ID: 701770}
\textbf{Source problem code:} catalogue entry
\subsection{Definitions and mathematical statement}
All groups considered are finite; no solubility hypothesis is imposed. A formation is a nonempty class of finite groups closed under isomorphisms, homomorphic images, and finite subdirect products. A formation \(\mathfrak F\) is hereditary if every subgroup of each member belongs to \(\mathfrak F\). It is saturated if, for every finite group \(X\),
\[
X/\Phi(X)\in\mathfrak F\quad\Longrightarrow\quad X\in\mathfrak F,
\]
where \(\Phi(X)\) is the intersection of the maximal subgroups of \(X\), with \(\Phi(1)=1\).

Write
\[
\pi(X)=\{p:p\text{ is prime and }p\mid |X|\},\qquad
\pi(\mathfrak F)=\bigcup_{X\in\mathfrak F}\pi(X).
\]
For \(H\leq X\), its core in \(X\) is \(\operatorname{Core}_X(H)=\bigcap_{x\in X}x^{-1}Hx\). Under the conventional formation-theoretic interpretation used in this draft, \(H\) is \(\mathfrak F\)-subnormal in \(X\) if either \(H=X\), or there is a chain
\[
H=H_0<H_1<\cdots<H_n=X
\]
in which \(H_{i-1}\) is maximal in \(H_i\) and
\[
H_i/\operatorname{Core}_{H_i}(H_{i-1})\in\mathfrak F
\qquad(1\leq i\leq n).
\]
The supplied record does not define \(\mathfrak F\)-subnormality, so this convention requires confirmation against the cited source.

Let \(\mathfrak F\) be a hereditary saturated formation. For \(P\leq G\), let \(N_G(P)=\{g\in G:g^{-1}Pg=P\}\), and let \(\operatorname{Syl}_p(G)\) denote the set of Sylow \(p\)-subgroups of \(G\). Define
\[
\mathrm{w}^{*}\mathfrak F=
\left\{G:\begin{array}{l}
G\text{ is finite},\quad \pi(G)\subseteq\pi(\mathfrak F),\\
N_G(P)\text{ is }\mathfrak F\text{-subnormal in }G\\
\text{for every }p\in\pi(G)\text{ and every }P\in\operatorname{Syl}_p(G)
\end{array}\right\}.
\]
The source states that \(\mathrm{w}^{*}\mathfrak F\) is a formation. The question is whether it must be saturated: for every hereditary saturated formation \(\mathfrak F\) and every finite group \(G\), must
\[
G/\Phi(G)\in\mathrm{w}^{*}\mathfrak F
\quad\Longrightarrow\quad
G\in\mathrm{w}^{*}\mathfrak F
\]
hold?
\subsection{Short English statement}
Given a hereditary saturated formation \(\mathfrak F\), consider the finite groups whose prime divisors occur in \(\mathfrak F\) and whose Sylow subgroup normalizers are \(\mathfrak F\)-subnormal. The source states that these groups form a formation and asks whether that formation is always saturated. The intended convention for \(\mathfrak F\)-subnormality requires source confirmation.
\subsection{Sources}
\begin{itemize}
\item Local UnsolvedMath registry, permanent ID \(701770\), source code \texttt{catalogue entry}, title \emph{Kourovka Notebook Problem 20.112}. The supplied registry status is \texttt{open}; this label is not an independently verified mathematical assessment.
\item \emph{The Kourovka Notebook (Unsolved Problems in Group Theory)}, problem 20.112, attributed in the supplied record to A. F. Vasil'ev and T. I. Vasil'eva, proposed in 2022. The background identifies the 20th issue (2022) and the 21st issue edition (2026), with identifier \texttt{arXiv:1401.0300v48}. These bibliographic details are supplied attribution and have not been independently checked.
\item The supplied collection description names E. I. Khukhro and V. D. Mazurov as editors of the Kourovka Notebook collection covering Issues 1--21 (1965--2026).
\end{itemize}
\end{document}
